% ============================================================================
%  Part X: Migration Guide
% ============================================================================
\section{Migration Guide}
\label{sec:migration}

The recommended path from an existing identifier format (UUIDv4, UUIDv7, ULID, KSUID, Snowflake, BIGSERIAL) to PQID is a five-phase rolling migration that preserves backward compatibility throughout.

\subsection{The five phases}

\textbf{Phase 1 --- Add the column.} Introduce PQID as a nullable column or field next to the existing identifier; both coexist, the existing identifier remains the primary key.

\begin{lstlisting}[language=SQL,caption={Adding a PQID column}]
ALTER TABLE events ADD COLUMN pqid CHAR(55);  -- B64URL, 55 chars
                                              -- or BYTEA(41) for binary
CREATE INDEX events_pqid_idx ON events (pqid);
\end{lstlisting}

\textbf{Phase 2 --- Update writers.} All writers populate PQID alongside the existing identifier on every new row. If the existing identifier carries a timestamp (UUIDv7, ULID, KSUID), derive PQID's timestamp from it; otherwise use current wall-clock time. The PQID's random tail must always be fresh from the CSPRNG --- deriving it from the old identifier preserves only the old identifier's entropy (e.g., 122 bits for UUIDv4), not 256.

\textbf{Phase 3 --- Backfill historical rows.} For rows whose PQID is null, generate synthesized PQIDs in a background batch job. If the existing identifier is timestamped, decode and reuse that timestamp; otherwise use a \code{created\_at} column, an autoincrement-to-time mapping, or a coarse epoch. The random tail is fresh per row, so backfilled PQIDs are guaranteed unique even when synthesized timestamps collide. Run in batches, throttled to avoid impacting production.

\textbf{Phase 4 --- Migrate readers.} Update read paths to accept either the old identifier or the PQID (length and format disambiguate). Update write paths to return PQID as the canonical identifier in API responses, while still accepting old identifiers as input. Migrate clients over time. Once all clients have migrated, the old identifier becomes accept-only. This phase is rate-limited by client deployment cadence; plan for months.

\textbf{Phase 5 --- Drop the old column.} Once no client or downstream system references the old identifier, drop it. This phase is irreversible without a restore-from-backup; verify no consumer references the old identifier by running it deprecated-but-present for at least one full cycle of all known consumers first.

\begin{lstlisting}[language=SQL,caption={Dropping the old identifier column}]
ALTER TABLE events DROP CONSTRAINT events_pkey;
ALTER TABLE events ADD PRIMARY KEY (pqid);
ALTER TABLE events DROP COLUMN old_id;
\end{lstlisting}

\subsection{Pitfalls}

\begin{itemize}
  \item \textbf{Don't derive PQID's random tail from the old identifier} --- the post-quantum margin lives in 256 fresh CSPRNG bits, and deriving from a 122-bit UUIDv4 caps the effective margin at the old format's entropy.
  \item \textbf{Don't treat the PQID timestamp as authoritative for ``when this row was created''} --- backfilled rows have synthesized timestamps. Use a separate \code{created\_at} column for that.
  \item \textbf{Don't mix Standard and HLC PQIDs in a single ordered stream} (Part \ref{sec:distributed}) --- pick one variant per table.
  \item \textbf{Validate the variant byte} on decode --- reject anything other than \code{0x01} or \code{0x02}.
  \item \textbf{Accept both encodings on input} --- some surfaces serialize as Crockford rather than Base64URL.
\end{itemize}

\subsection{Cross-region replication during migration}

If the system replicates rows across regions, complete Phase 1 (add column) and Phase 3 (backfill) globally before starting Phase 2 (writer updates). Otherwise replicas without the column receive populated-PQID writes from peers that have it, leading to mixed states the migration plan must handle anyway.

\subsection{Legacy systems with no clean migration path}

When the old identifier is exposed in URLs that customers have bookmarked or in third-party integrations whose update cadence is unknown, keep the old identifier permanently and use PQID as a parallel identifier for new internal use cases. Both coexist indefinitely.
